% Original: PDF strana 22, donji slajd.
\slajdid{02-s044}
\begin{frame}[label=02-s044]{Primer: minimalni proizvod zbirova}
\setlength{\parskip}{0pt}
% element: 02-s044:e01
Za minimalnu funkciju u obliku proizvoda zbirova prekrivamo sve logičke nule površinama najvišeg reda, koristeći po potrebi i neodređeno stanje \(b\).
\par\vspace{2mm}
\begin{columns}[T,onlytextwidth]
\begin{column}{.39\textwidth}\centering
% element: 02-s044:e02
% Original prikazuje samo grupu celog reda DC=11; ostale grupe se ne dopisuju.
\begingroup
\fontsize{12}{13.3}\selectfont
\renewcommand{\small}{\fontsize{12}{13.3}\selectfont}
\scalebox{.75}{%
\begin{karnaugh-map}(label=middle,implicantcolors={DEcrvena,DEcrvena,DEcrvena})[4][4][1][$A$][$B$][$C$][$D$]
\minterms{0,4,5,8,10}
\terms{1}{b}
\maxterms{2,3,6,7,9,11,12,13,14,15}
\implicant{12}{14}
\end{karnaugh-map}}
\endgroup

\end{column}
\begin{column}{.57\textwidth}
% element: 02-s044:e03
Površine objedinjuju kombinacije za koje je moguće sažimanje izraza.
\par\vspace{2mm}
U zbirovima ostaju \alert{\textbf{komplementi literala}} koji se \alert{\textbf{ne menjaju}} na površini.
\par\vspace{5mm}
% element: 02-s044:e04
\par\Rezultat{\(F=(D+\bar B)(\bar D+\bar C)(\bar D+\bar A)\)}
\end{column}
\end{columns}
\end{frame}
